\subsection{Модуль машинного обучения (Analytics-ML)}

Наряду с серверной частью и клиентским приложением, система включает модуль машинного обучения, отвечающий за предиктивную аналитику. Модуль реализован на Python с использованием FastAPI и предоставляет API для прогнозирования и обнаружения аномалий. На рисунке~\ref{ml-arch} представлена архитектура модуля машинного обучения.

\begin{figure}
  \centering
  \includegraphics[width=0.95\textwidth,height=0.85\textheight,keepaspectratio]{analytics-ml.png}
  \caption{Архитектура модуля машинного обучения}
  \label{ml-arch}
\end{figure}

\subsubsection{Выбор алгоритмов}

Все задачи предиктивной аналитики в системе формулируются на табулярных данных, образованных агрегированными показателями животных (надои, активность, жвачка, SCC) за скользящие временные окна. Для таких данных методы градиентного бустинга показывают состояние-of-the-art результаты~\cite{XGBoost}, превосходя нейросетевые подходы на выборках умеренного размера (сотни---тысячи записей), характерных для одной фермы.

Система автоматически выбирает наилучший алгоритм градиентного бустинга среди трёх реализаций: XGBoost~\cite{XGBoost}, LightGBM~\cite{LightGBM} и CatBoost~\cite{CatBoost}. Выбор обусловлен различиями в подходах к построению деревьев: XGBoost использует уровень-wise рост, обеспечивая стабильность на сбалансированных данных; LightGBM применяет лист-wise рост с гистограммной оптимизацией, что ускоряет обучение на больших выборках; CatBoost использует ordered boosting, устойчивый к утечке целевой переменной (target leakage). Гипероптимизация Optuna проводит испытания равномерно по всем трём бэкендам и выбирает лучший по результатам кросс-валидации.

Для кластеризации стада используется алгоритм KMeans~\cite{KMeans}, либо HDBSCAN~\cite{HDBSCAN} при наличии кластеров произвольной формы; выбор производится автоматически по метрике силуэта. Для обнаружения аномалий оборудования применяется Isolation Forest~\cite{IsolationForest}, не требующий размеченных данных и устойчивый к высокой доле аномалий. Для прогнозирования на уровне стада используется Prophet~\cite{Prophet}, автоматически декомпозирующий временной ряд на тренд, сезонность и эффекты особых дней.

Разработано 11 моделей машинного обучения (таблица~\ref{mlmodels}).

\begin{table}
  \caption{Модели машинного обучения}
  \small
  \begin{tabularx}{\textwidth}{|p{3.2cm}|p{2.8cm}|l|X|}
    \hline
    \textbf{Модель} & \textbf{Алгоритм} & \textbf{Метрика} & \textbf{Назначение} \\
    \hline
    Прогноз надоев (корова) & XGBoost, квантильная регрессия & MAE & Дневной надой с q10/q90 интервалами \\
    \hline
    Прогноз надоев (стадо) & Prophet & --- & Тренд надоя стада с сезонностью \\
    \hline
    Риск мастита & XGBoost / LightGBM / CatBoost & ROC-AUC & Вероятность мастита \\
    \hline
    Предупреждение кетоза & XGBoost / LightGBM / CatBoost & ROC-AUC & Риск метаболического нарушения \\
    \hline
    Обнаружение охоты & XGBoost / LightGBM / CatBoost & ROC-AUC & Состояние (в охоте / нет) \\
    \hline
    Риск выбраковки & XGBoost / LightGBM / CatBoost & $R^2$ & Дни до выбраковки \\
    \hline
    Оценка BCS & XGBoost / LightGBM / CatBoost & MAE & Упитанность (1,0--5,0) \\
    \hline
    Кластеризация стада & KMeans / HDBSCAN & Silhouette & Сегментация стада \\
    \hline
    Аномалии оборудования & Isolation Forest & Доля аномалий & Аномалии работы роботов \\
    \hline
    Рекомендации по корму & XGBoost / LightGBM / CatBoost & MAE & Объём корма \\
    \hline
    Комплексное здоровье & MultiOutput\newline Classifier & ROC-AUC & Мастит + кетоз + охота совместно \\
    \hline
  \end{tabularx}
  \label{mlmodels}
\end{table}

\subsubsection{Обоснование выбора метрик}

Для задач классификации (мастит, кетоз, охота) выбрана метрика ROC-AUC (Area Under the Receiver Operating Characteristic Curve), поскольку она инвариантна к порогу принятия решения и соотношению классов. Это критично для veterinary-задач, где положительный класс составляет менее 5~процентов наблюдений: метрики accuracy и precision при таком дисбалансе неинформативны. ROC-AUC оценивает способность модели ранжировать наблюдения по степени риска, что соответствует прикладной задаче --- выделить животных с наибольшей вероятностью заболевания для приоритетного осмотра.

Для задач регрессии (прогноз надоев, рекомендации по корму, BCS) выбрана MAE (Mean Absolute Error), а не MSE/ RMSE. Обоснование: MAE менее чувствительна к выбросам, которые характерны для данных доильных роботов (пропущенное доение, сбой датчика). Кроме того, MAE измеряется в тех же единицах, что и целевая переменная (кг молока, кг корма), что упрощает интерпретацию результатов фермером.

Для кластеризации выбрана метрика силуэта, оценивающая компактность и разделимость кластеров без необходимости внешних меток.

Оценка всех моделей производится методом $k$-кратной кросс-валидации, где число фолдов адаптируется к объёму данных: $k = \min(5, \max(2, N / n))$, где $N$ --- количество наблюдений, $n$ --- минимальный размер фолда.

\subsubsection{Формирование обучающих меток}

Особенностью задач обнаружения заболеваний является отсутствие надёжных размеченных данных: ветеринарные записи неполны, а диагноз часто ставится с задержкой. Для решения этой проблемы применён подход обучения на положительных и неразмеченных данных (PU-learning)~\cite{PULearning}.

Метка <<мастит>> формируется на основе комбинации клинически обоснованных правил: SCC~>~300\,000 клеток/мл; либо SCC~>~200\,000 с трендом роста~>~1,5; либо SCC~>~150\,000 при повышенной электропроводности~>~55 мСм/см; либо отклонение надоя~>~15~процентов при SCC~>~100\,000. Аналогичные правила на основе FPR и жвачки используются для кетоза, а на основе соотношения активности и жвачки --- для охоты.

При наличии подтверждённых ветеринарных диагнозов они замещают эвристические метки с увеличением веса наблюдения в 3 раза, что позволяет модели постепенно переходить от semi-supervised к fully-supervised режиму по мере накопления данных.

\subsubsection{Прогноз надоев с доверительными интервалами}

Для прогнозирования надоев на корову обучаются три модели квантильной регрессии: для медианы (q50), нижней границы (q10) и верхней границы (q90). Объективная функция Quantile Loss для квантиля $\alpha$:

$$L_\alpha(y, \hat{y}) = \begin{cases} \alpha \cdot (y - \hat{y}), & \text{если } y \geq \hat{y} \\ (1 - \alpha) \cdot (\hat{y} - y), & \text{если } y < \hat{y} \end{cases}$$

где $y$ --- фактическое значение, $\hat{y}$ --- прогноз. Интервал q10--q90 содержит 80~процентов распределения прогнозов, позволяя фермеру количественно оценивать неопределённость.

\subsubsection{Инженерия признаков}

Для каждой модели формируется расширенный набор признаков на основе данных из PostgreSQL с использованием LATERAL JOIN для вычисления скользящих агрегатов. Признаки включают:

\begin{enumerate}
    \item лаговые --- значения целевого показателя за предыдущие 1, 7, 14, 30 дней;
    \item скользящие статистики --- среднее, стандартное отклонение, коэффициент вариации за окна 7, 14, 30 дней;
    \item трендовые --- отношение recent/baseline, наклон линейной регрессии за 7 и 14 дней;
    \item производные --- FPR, соотношение надой/корм, соотношение активность/жвачка;
    \item физиологические --- DIM, номер лактации, возраст;
    \item метеорологические --- температура, влажность, индекс температурно-влажностного стресса (THI);
    \item исторические --- количество ветеринарных обработок за 180 дней, дней с последней обработки;
    \item временные ряды --- энтропия, автокорреляция на лаге 1 и 7, сила тренда и сезонности.
\end{enumerate}

Для модели прогноза надоев дополнительно вычисляются TSFRESH-подобные признаки: приближённая энтропия, коэффициент эксцесса, максимальная просадка (drawdown), сила сезонности по STL-декомпозиции. Эти признаки позволяют модели различать стабильных и нестабильных коров.

Обработка пропущенных значений выполняется медианой по стаду, вычисленной на этапе обучения и сохраняемой вместе с моделью для обеспечения консистентности при прогнозировании. Дополнительно выполняется отбор признаков с помощью RandomForest importance: отбираются признаки с важностью выше медианной.

\subsubsection{MLOps: оптимизация, мониторинг дрейфа и управление моделями}

Оптимизация гиперпараметров осуществляется библиотекой Optuna~\cite{Optuna} с TPE-сэмплером (Tree-structured Parzen Estimator) в 8-мерном пространстве (learning\_rate, max\_depth, n\_estimators, min\_child\_weight, subsample, colsample\_bytree, reg\_alpha, reg\_lambda). Каждый испытание оценивается кросс-валидацией; по умолчанию выполняется 30 испытаний с таймаутом 120 секунд.

Мониторинг дрейфа реализован тремя методами. \textbf{Prediction drift} --- Z-тест: если $z = |\bar{x}_{\text{recent}} - \bar{x}_{\text{baseline}}| / \sigma_{\text{baseline}} > 2{,}0$, фиксируется сдвиг распределения прогнозов. \textbf{Feature drift} --- для каждого признака вычисляется Z-тест и критерий Колмогорова--Смирнова; дрейф фиксируется при $z > 2{,}5$ или $p_{\text{KS}} < 0{,}05$. \textbf{MMD-drift} --- Maximum Mean Discrepancy с RBF-ядром, оценивающий расстояние между распределениями в пространстве признаков; порог MMD~>~0{,}1.

При обнаружении дрейфа модель помечается для скорого переобучения, которое выполняется через APScheduler с настраиваемым расписанием (по умолчанию еженедельно).

Для управления моделями реализован подход <<чемпион--претендент>>: переобученная модель (претендент) сравнивается с текущей (чемпионом) на тех же данных; модель повышается только при улучшении метрики более чем на 1~процент. Это предотвращает деградацию при переобучении на зашумлённых данных.

Объяснимость прогнозов реализована через SHAP~\cite{SHAP} с использованием TreeExplainer для древовидных моделей. Отслеживание экспериментов осуществляется через MLflow~\cite{MLflow}. Для ускорения инференса модели экспортируются в формат ONNX.
